\documentclass[10pt,a4paper]{amsart}
\usepackage{amsmath,amssymb,amsthm,graphicx}
% Retained minimal preamble for the Stacks Project chapter more-algebra.tex (More on Algebra).
% The chapter file pinned-original-source.tex carries only \input{preamble} on its first line, so no
% preamble text lives inside the chapter itself.  The Stacks Project uses one single shared preamble for
% every chapter; the snapshot stages it identically as sources/stacks-chow/preamble.tex and
% sources/stacks-etale-cohomology/preamble.tex.  The block below is copied mechanically, line for line,
% from sources/stacks-chow/preamble.tex lines 14-18 (the Xy-pic setup) and 154-217 (the native theorem
% declarations together with the mathematical macros).  Nothing was added, removed or redefined; the
% wrapper class and the amsmath/amssymb/amsthm/graphicx base packages are supplied by the assembler.
% Chapter: more-algebra.tex -- More on Algebra.
% For commutative diagrams we use Xy-pic
\usepackage[all]{xy}

% We use 2cell for 2-commutative diagrams.
\xyoption{2cell}

% Theorem environments.
%
\theoremstyle{plain}
\newtheorem{theorem}[subsection]{Theorem}
\newtheorem{proposition}[subsection]{Proposition}
\newtheorem{lemma}[subsection]{Lemma}

\theoremstyle{definition}
\newtheorem{definition}[subsection]{Definition}
\newtheorem{example}[subsection]{Example}
\newtheorem{exercise}[subsection]{Exercise}
\newtheorem{situation}[subsection]{Situation}

\theoremstyle{remark}
\newtheorem{remark}[subsection]{Remark}
\newtheorem{remarks}[subsection]{Remarks}

\numberwithin{equation}{subsection}

% Macros
%
\def\lim{\mathop{\mathrm{lim}}\nolimits}
\def\colim{\mathop{\mathrm{colim}}\nolimits}
\def\Spec{\mathop{\mathrm{Spec}}}
\def\Hom{\mathop{\mathrm{Hom}}\nolimits}
\def\Ext{\mathop{\mathrm{Ext}}\nolimits}
\def\SheafHom{\mathop{\mathcal{H}\!\mathit{om}}\nolimits}
\def\SheafExt{\mathop{\mathcal{E}\!\mathit{xt}}\nolimits}
\def\Sch{\mathit{Sch}}
\def\Mor{\mathop{\mathrm{Mor}}\nolimits}
\def\Ob{\mathop{\mathrm{Ob}}\nolimits}
\def\Sh{\mathop{\mathit{Sh}}\nolimits}
\def\NL{\mathop{N\!L}\nolimits}
\def\CH{\mathop{\mathrm{CH}}\nolimits}
\def\proetale{{pro\text{-}\acute{e}tale}}
\def\etale{{\acute{e}tale}}
\def\QCoh{\mathit{QCoh}}
\def\Ker{\mathop{\mathrm{Ker}}}
\def\Im{\mathop{\mathrm{Im}}}
\def\Coker{\mathop{\mathrm{Coker}}}
\def\Coim{\mathop{\mathrm{Coim}}}

% Boxtimes
%
\DeclareMathSymbol{\boxtimes}{\mathbin}{AMSa}{"02}

%
% Macros for moduli stacks/spaces
%
\def\QCohstack{\mathcal{QC}\!\mathit{oh}}
\def\Cohstack{\mathcal{C}\!\mathit{oh}}
\def\Spacesstack{\mathcal{S}\!\mathit{paces}}
\def\Quotfunctor{\mathrm{Quot}}
\def\Hilbfunctor{\mathrm{Hilb}}
\def\Curvesstack{\mathcal{C}\!\mathit{urves}}
\def\Polarizedstack{\mathcal{P}\!\mathit{olarized}}
\def\Complexesstack{\mathcal{C}\!\mathit{omplexes}}
% \Pic is the operator that assigns to X its picard group, usage \Pic(X)
% \Picardstack_{X/B} denotes the Picard stack of X over B
% \Picardfunctor_{X/B} denotes the Picard functor of X over B
\def\Pic{\mathop{\mathrm{Pic}}\nolimits}
\def\Picardstack{\mathcal{P}\!\mathit{ic}}
\def\Picardfunctor{\mathrm{Pic}}
\def\Deformationcategory{\mathcal{D}\!\mathit{ef}}
\usepackage{cleveref}
\crefname{originallemma}{original lemma}{original lemmas}
\Crefname{originallemma}{Original lemma}{Original lemmas}
\makeatletter
\newlabel{lemma-base-change-invariants}{{0BRH}{source}{}{sourcepointer-1}{}}
\newlabel{lemma-base-change-invariants@cref}{{[originallemma][0][]0BRH}{[1][0][]source}{}{}{}}
\makeatother
\crefname{originallemma}{original lemma}{original lemmas}
\Crefname{originallemma}{Original lemma}{Original lemmas}
\makeatletter
\newlabel{lemma-one-orbit}{{0BRI}{source}{}{sourcepointer-2}{}}
\newlabel{lemma-one-orbit@cref}{{[originallemma][0][]0BRI}{[1][0][]source}{}{}{}}
\makeatother
\crefname{originallemma}{original lemma}{original lemmas}
\Crefname{originallemma}{Original lemma}{Original lemmas}
\makeatletter
\newlabel{algebra-lemma-integral-no-inclusion}{{00GT}{source}{}{sourcepointer-3}{}}
\newlabel{algebra-lemma-integral-no-inclusion@cref}{{[originallemma][0][]00GT}{[1][0][]source}{}{}{}}
\makeatother
\crefname{originallemma}{original lemma}{original lemmas}
\Crefname{originallemma}{Original lemma}{Original lemmas}
\makeatletter
\newlabel{algebra-lemma-integral-going-up}{{00GU}{source}{}{sourcepointer-4}{}}
\newlabel{algebra-lemma-integral-going-up@cref}{{[originallemma][0][]00GU}{[1][0][]source}{}{}{}}
\makeatother
\crefname{originallemma}{original lemma}{original lemmas}
\Crefname{originallemma}{Original lemma}{Original lemmas}
\makeatletter
\newlabel{algebra-lemma-chinese-remainder}{{00DT}{source}{}{sourcepointer-5}{}}
\newlabel{algebra-lemma-chinese-remainder@cref}{{[originallemma][0][]00DT}{[1][0][]source}{}{}{}}
\makeatother
\crefname{originallemma}{original lemma}{original lemmas}
\Crefname{originallemma}{Original lemma}{Original lemmas}
\makeatletter
\newlabel{fields-lemma-normal-case}{{030M}{source}{}{sourcepointer-6}{}}
\newlabel{fields-lemma-normal-case@cref}{{[originallemma][0][]030M}{[1][0][]source}{}{}{}}
\makeatother
\crefname{originallemma}{original lemma}{original lemmas}
\Crefname{originallemma}{Original lemma}{Original lemmas}
\makeatletter
\newlabel{fields-lemma-primitive-element}{{030N}{source}{}{sourcepointer-7}{}}
\newlabel{fields-lemma-primitive-element@cref}{{[originallemma][0][]030N}{[1][0][]source}{}{}{}}
\makeatother
\title{Residue field extensions of a finite group action are normal, and the decomposition group surjects onto the automorphism group}
\author{The Stacks Project Authors}
\date{stacks-project:more-algebra.tex:local-snapshot-20261012; GNU FDL 1.2 or later}
\begin{document}
\maketitle
\setcounter{subsection}{1}
\noindent\textit{Editorial note.} All delivered mathematical text is the author's own, copied from the pinned Stacks Project chapter file more-algebra.tex; no statement or proof step was rewritten, reordered or removed. The counted statement is the labeled lemma environment at source lines 32703-32716 (canonical tag 0BRJ) and the counted proof is the proof environment at source lines 32718-32774, which immediately follows it and is its own proof.\par\medskip
\noindent\textit{Editorial note.} The statement carries one footnote recalling that the notation Gal is used only for Galois extensions; the footnote is part of the author's own statement inside the counted span.\par\medskip
\noindent\textit{Editorial note.} The proof refers twice to results of this chapter and five times to results of other chapters, namely integrality statements, the Chinese remainder theorem, the splitting of normal extensions and the primitive element theorem; all seven are navigation stubs with their roles recorded in the reasons.\par\medskip
\noindent\textit{Editorial note.} The chapter file pinned-original-source.tex carries only an input of the shared preamble on its first line, so no preamble text lives inside the chapter. The retained file sources/stacks-more-algebra/retained-preamble.tex is a mechanical, line-for-line copy of the project's single shared preamble as staged in the snapshot (Xy-pic setup plus the native theorem declarations and mathematical macros). Nothing was added, removed or redefined, and the excerpt serves only as preamble, never as a span coordinate.\par\medskip
\noindent\textit{Editorial note.} The source URL is the source-level Stacks Project URL, as required when more than one card is produced from one source; the per-claim tag is carried by claim\_key and by this note. The counted statement and proof cite no bibliography entry, so this card needs no bibliography.\par\medskip
\noindent\textit{Allowed standard prerequisites (editorial).}\begin{itemize}
\item The invariant subring of a finite group action is integrally closed in the ambient ring, so primes lying over a given prime exist and the maximal ideals over a maximal ideal form one orbit. used at proof lines 32726-32733. Restricts the analysis to the finitely many primes over the chosen prime, which form a single orbit under the group.
\item For a finite product of quotients by the primes in one orbit the Chinese remainder theorem makes the diagonal map surjective. invoked at proof lines 32734-32736. Allows an element vanishing at all but one of the primes to be chosen, which is the input to the normality argument.
\item Every finite separable extension admits a primitive element, and any normal extension splits as the tensor product of a purely inseparable and a Galois part. invoked at proof lines 32754-32762. Splits the residue field extension into its purely inseparable and Galois parts and produces the element whose degree is compared with the order of the group.
\end{itemize}\medskip
\ifdefined\hypertarget\hypertarget{sourcepointer-1}{}\fi
\noindent\textit{Source reference.} Original lemma 0BRH: sources/stacks-more-algebra/pinned-original-source.tex, source lines 32656--32656. The reference is invoked at proof lines 32719-32720 to see that the localized invariant ring is the invariant ring of the localized ring; the reduction that uses it is the first step inside the counted span.\par\medskip
\ifdefined\hypertarget\hypertarget{sourcepointer-2}{}\fi
\noindent\textit{Source reference.} Original lemma 0BRI: sources/stacks-more-algebra/pinned-original-source.tex, source lines 32680--32680. The reference supplies the finiteness and single orbit of the primes over the chosen maximal prime, used at proof lines 32731-32733; the Chinese remainder and polynomial arguments built on it are inside the counted span.\par\medskip
\ifdefined\hypertarget\hypertarget{sourcepointer-3}{}\fi
\noindent\textit{Source reference.} Original lemma 00GT: sources/stacks-chow/tags.txt, source lines 619--619. The reference is one of the two results cited at proof lines 32728-32729 to identify the maximal ideals over the maximal prime; it is a standard pointer and its use is inside the counted span.\par\medskip
\ifdefined\hypertarget\hypertarget{sourcepointer-4}{}\fi
\noindent\textit{Source reference.} Original lemma 00GU: sources/stacks-chow/tags.txt, source lines 620--620. The reference is the second of the two results cited at proof lines 32728-32729 for the same identification; it is a standard pointer, and everything built on it stays inside the counted span.\par\medskip
\ifdefined\hypertarget\hypertarget{sourcepointer-5}{}\fi
\noindent\textit{Source reference.} Original lemma 00DT: sources/stacks-chow/tags.txt, source lines 509--509. The reference makes the diagonal map to the product of the quotients surjective, used at proof lines 32734-32736 to choose an element vanishing outside one prime; the polynomial argument consuming that choice is inside the counted span.\par\medskip
\ifdefined\hypertarget\hypertarget{sourcepointer-6}{}\fi
\noindent\textit{Source reference.} Original lemma 030M: sources/stacks-chow/tags.txt, source lines 3953--3953. The reference splits the normal residue field extension into a purely inseparable and a Galois part, used at proof lines 32754-32758; the surjectivity argument that follows is inside the counted span.\par\medskip
\ifdefined\hypertarget\hypertarget{sourcepointer-7}{}\fi
\noindent\textit{Source reference.} Original lemma 030N: sources/stacks-chow/tags.txt, source lines 3954--3954. The reference supplies a primitive element for the Galois part and its infinite case analogue, used at proof lines 32759-32762; the comparison of its degree with the order of the group is inside the counted span.\par\medskip
\begin{lemma}
\label{lemma-one-orbit-geometric}
Let $G$ be a finite group acting on a ring $R$. Let $\mathfrak q \subset R$
be a prime lying over $\mathfrak p \subset R^G$. Then
$\kappa(\mathfrak q)/\kappa(\mathfrak p)$ is an algebraic normal
extension and the map
$$
D = \{\sigma \in G \mid \sigma(\mathfrak q) = \mathfrak q\}
\longrightarrow
\text{Aut}(\kappa(\mathfrak q)/\kappa(\mathfrak p))
$$
is surjective\footnote{Recall that we use the notation $\text{Gal}$
only in the case of Galois extensions.}.
\end{lemma}

\begin{proof}
With $A = (R^G)_\mathfrak p$ and $B = A \otimes_{R^G} R$ we see that $A = B^G$
as localization is flat, see Lemma \ref{lemma-base-change-invariants}.
Observe that $\mathfrak pA$ and $\mathfrak qB$ are prime ideals,
$D$ is the stabilizer of $\mathfrak qB$, and
$\kappa(\mathfrak p) = \kappa(\mathfrak pA)$ and
$\kappa(\mathfrak q) = \kappa(\mathfrak qB)$. Thus we may replace
$R$ by $B$ and assume that $\mathfrak p$ is a maximal ideal.
Since $R^G \subset R$ is an integral ring extension, we find
that the maximal ideals of $R$ are exactly the primes lying over
$\mathfrak p$ (follows from
Algebra, Lemmas \ref{algebra-lemma-integral-no-inclusion} and
\ref{algebra-lemma-integral-going-up}).
By Lemma \ref{lemma-one-orbit} there are finitely many of them
$\mathfrak q = \mathfrak q_1, \mathfrak q_2, \ldots, \mathfrak q_m$
and they form a single orbit for $G$.
By the Chinese remainder theorem
(Algebra, Lemma \ref{algebra-lemma-chinese-remainder}) the map
$R \to \prod_{j = 1, \ldots, m} R/\mathfrak q_j$ is surjective.

\medskip\noindent
First we prove that the extension is normal. Pick an element
$\alpha \in \kappa(\mathfrak q)$. We have to show that the
minimal polynomial $P$ of $\alpha$ over $\kappa(\mathfrak p)$
splits completely. By the above we can choose
$a \in \mathfrak q_2 \cap \ldots \cap \mathfrak q_m$
mapping to $\alpha$ in $\kappa(\mathfrak q)$.
Consider the polynomial $Q = \prod_{\sigma \in G} (T - \sigma(a))$
in $R^G[T]$. The image of $Q$ in $R[T]$ splits completely
into linear factors, hence the same is true for its
image in $\kappa(\mathfrak q)[T]$. Since $P$ divides
the image of $Q$ in $\kappa(\mathfrak p)[T]$ we conclude
that $P$ splits completely into linear factors over
$\kappa(\mathfrak q)$ as desired.

\medskip\noindent
Since $\kappa(\mathfrak q)/\kappa(\mathfrak p)$ is normal we may assume
$\kappa(\mathfrak q) = \kappa_1 \otimes_{\kappa(\mathfrak p)} \kappa_2$
with $\kappa_1/\kappa(\mathfrak p)$ purely inseparable and
$\kappa_2/\kappa(\mathfrak p)$ Galois, see
Fields, Lemma \ref{fields-lemma-normal-case}. Pick
$\alpha \in \kappa_2$ which generates $\kappa_2$ over $\kappa(\mathfrak p)$
if it is finite and a subfield of degree $> |G|$ if it is infinite
(to get a contradiction).
This is possible by Fields, Lemma \ref{fields-lemma-primitive-element}.
Pick $a$, $P$, and $Q$ as in the previous paragraph.
If $\alpha' \in \kappa_2$ is a Galois conjugate of $\alpha$ over
$\kappa(\mathfrak p)$, then the fact that $P$ divides the
image of $P$ in $\kappa(\mathfrak p)[T]$ shows there exists a
$\sigma \in G$ such that $\sigma(a)$ maps to $\alpha'$. By our choice of
$a$ (vanishing at other maximal ideals) this implies $\sigma \in D$ and
that the image of $\sigma$ in
$\text{Aut}(\kappa(\mathfrak q)/\kappa(\mathfrak p))$
maps $\alpha$ to $\alpha'$. Hence the surjectivity or the
desired absurdity
in case $\alpha$ has degree $> |G|$ over $\kappa(\mathfrak p)$.
\end{proof}

\end{document}
